\chapter{A numerical two-trait example from beginning to end}

\section{Specify \(\G\) and two curvature matrices}

Let
\begin{equation}
\G=
\begin{bmatrix}
2&0\\0&0.5
\end{bmatrix},
\end{equation}
so
\begin{equation}
\G^{1/2}=
\begin{bmatrix}
\sqrt2&0\\0&1/\sqrt2
\end{bmatrix}.
\end{equation}
Consider two populations with
\begin{equation}
\Hh_1=
\begin{bmatrix}
1&0.3\\0.3&1.4
\end{bmatrix},
\qquad
\Hh_2=
\begin{bmatrix}
1&-0.3\\-0.3&1.4
\end{bmatrix}.
\end{equation}
Both are positive definite because their determinants are
\begin{equation}
1(1.4)-0.3^2=1.31>0
\end{equation}
and their first diagonal entry is positive.

\section{Construct the channel operators}

Using \cref{eq:S2D},
\begin{equation}
\Sop_1=
\begin{bmatrix}
2&0.3\\0.3&0.7
\end{bmatrix},
\qquad
\Sop_2=
\begin{bmatrix}
2&-0.3\\-0.3&0.7
\end{bmatrix}.
\end{equation}
The eigenvalues of either matrix are
\begin{align}
\lambda_{1,2}
&=\frac{2+0.7}{2}
\pm\half\sqrt{(2-0.7)^2+4(0.3)^2}\\
&=1.35\pm\half\sqrt{2.05}\\
&\approx2.0659,\;0.6341.
\end{align}
The gap is
\begin{equation}
\delta\approx1.4318.
\end{equation}

For population 1,
\begin{equation}
\phi_1
=\half\operatorname{atan2}(0.6,1.3)
\approx12.39^\circ.
\end{equation}
For population 2,
\begin{equation}
\phi_2\approx-12.39^\circ.
\end{equation}
Thus the whitened channel angle is
\begin{equation}
\theta_{12}=|\phi_1-\phi_2|
\approx24.78^\circ.
\end{equation}

\section{Compute geometric discrepancy}

For \(k=1\),
\begin{equation}
d_C^2=I_V=\sin^2\theta_{12}.
\end{equation}
Numerically,
\begin{equation}
d_C^2\approx\sin^2(24.78^\circ)\approx0.176.
\end{equation}
The overlap is therefore about \(1-0.176=0.824\) on this squared projection scale.

\section{Map to phenotype-space directions}

From \cref{eq:phenotypeAngle2D},
\begin{equation}
\tan\alpha_i
=\sqrt{\frac{0.5}{2}}\tan\phi_i
=0.5\tan\phi_i.
\end{equation}
Therefore
\begin{equation}
\alpha_1\approx6.26^\circ,
\qquad
\alpha_2\approx-6.26^\circ.
\end{equation}
The ordinary phenotype-space angle is smaller because mapping back stretches the high-variance first axis more strongly.

\section{Calibrate a hypothetical ridge load}

Suppose an associated ridge model has
\begin{equation}
\Delta g=1.5,
\qquad
\Delta h=1.0.
\end{equation}
Then
\begin{align}
L_{\mathrm{cov}}
&=\frac{1.5\times1.0}{4}(0.176)\\
&\approx0.066.
\end{align}
The corresponding bounded log-scale isolation is
\begin{equation}
I_L=1-e^{-0.066}\approx0.064.
\end{equation}

\begin{cautionbox}
The number \(0.176\) is pure geometry. The number \(0.066\) is a fitness load only because \(\Delta g\) and \(\Delta h\) have supplied a biological scale. Different covariance and curvature contrasts would map the same angle to a different load.
\end{cautionbox}
